\subsection{Picard restriction at interfaces of reducible fibres}

\begin{lemma}[C.2.INT.1; Proved elliptic-image restriction]
\label{st:C.2.INT.1}
Let $Y$ be a smooth compact complex threefold with
$H^2(Y,\Q)=0$, let $S_i$ be an irreducible component of
a connected reducible divisor $S$, and let
$h:Z\to S_i\hookrightarrow Y$ be the minimal resolution
of its normalization. Suppose that the minimal model
$\sigma:Z\to Z_0$ is a Hopf surface or a non-K\"ahler
elliptic surface. Every strict transform $E\subset Z$ of
an elliptic curve on $Z_0$ has a one-dimensional image
under $h$, and the normalization of that image is elliptic.
The actual restriction homomorphism
$\Pic^0(Y)\to\Pic^0(E)$ is nonconstant for each such $E$.
In particular such a component cannot meet a component
whose minimal model is an Enoki surface.
\end{lemma}
\begin{proof}
Choose an adjacent component $S_j$. The actual pullback
of its canonical divisor section gives a nonzero effective
Cartier divisor
\[
 A=h^*S_j>0,\qquad \OO_Z(A)=h^*\OO_Y(S_j).
 \tag{C.INT.1}
\]
Its first Chern class is torsion. For some $q>0$,
$\OO_Y(qS_j)\in\Pic^0(Y)$, so the restriction homomorphism
$\rho:\Pic^0(Y)\to\Pic^0(Z)$ has in its image the line
$\OO_Z(qA)$. That line has strictly positive degree for
every Gauduchon metric on $Z$.

We spell out the blowdown comparison. Put $A_0=\sigma_*A$.
The difference $A-\sigma^*A_0$ is exceptional and has
zero real Chern class: both $A$ and its pushforward have
zero real class. The negative definite exceptional
intersection matrix makes every coefficient of this
difference zero. Thus
\[
 A=\sigma^*A_0,\qquad A_0>0.
 \tag{C.INT.2}
\]
Nonzeroness follows by the same negative definiteness if
$A_0$ were zero. Pullback is an isomorphism
$\Pic^0(Z_0)\simeq\Pic^0(Z)$; the integral blowup
formula and $R\sigma_*\OO_Z=\OO_{Z_0}$ identify both
the lattice and the coefficient space defining it.
We therefore regard $\rho$ as a map into $\Pic^0(Z_0)$.
Its image contains $\OO_{Z_0}(qA_0)$.

We first prove the stronger assertion that the composite
restriction to $E$ cannot be zero. Suppose that it is zero.
This situation would in particular arise if $h(E)$ were
a point or had rational normalization.
In the latter case the map from the normal curve $E$
factors as $E\to\PP^1\to Y$. In both cases every line
from $\Pic^0(Y)$ restricts trivially to $E$. Here the
pullback to $\PP^1$ is in $\Pic^0(\PP^1)=0$; this is
the specified map giving the vanishing. The restriction
of $\sigma$ to $E$ is an isomorphism onto its smooth
elliptic image $E_0$.

For a Hopf $Z_0$ the homomorphism
\[
 H^1(Z_0,\R)\longrightarrow H^1(E_0,\R)
 \tag{C.INT.3}
\]
is nonzero. On a primary Hopf surface this follows from
its universal cover $\C^2\setminus\{0\}$: a finite image
of $\pi_1(E_0)$ in $\pi_1(Z_0)=\Z$ would lift a finite
unramified cover of $E_0$ to that cover. Its two coordinate
functions would be constant, contradicting the nonconstant
inclusion of the curve. For a secondary Hopf surface,
pull back along a finite unramified primary cover. A
component above $E_0$ is an elliptic curve. Transfer makes
$H^1(Z_0,\R)\to H^1(Z'_0,\R)$ an isomorphism between
one-dimensional spaces, and the preceding primary
argument proves (C.INT.3) after this pullback, hence
before it. The real-to-holomorphic coefficient injection
of D.1.7, applied on $E_0$, shows that
$H^1(Z_0,\OO)\to H^1(E_0,\OO)$ is nonzero. Since
$\Pic^0(Z_0)=\C^*$, the identity component of the
kernel of $\Pic^0(Z_0)\to\Pic^0(E_0)$ is zero.
Connectedness of $\Pic^0(Y)$ would therefore make
$\rho$ zero, contradicting (C.INT.1)--(C.INT.2).

For a non-K\"ahler elliptic $p:Z_0\to B$, every curve
on its minimal model is a reduced elliptic fibre.
The actual covering construction of
\cite[Lemma~2.2 and \S2]{BrinzanescuMoraru2005}
provides a finite cyclic map $\psi:Z'_0\to Z_0$ and a
principal elliptic bundle $p':Z'_0\to B'$ such that a
reduced fibre $E'$ above $E_0$ maps to $E_0$ by a finite
isogeny. To include multiple fibres, take a base cover
whose branching orders are divisible by their
multiplicities and normalize the fibre product. In the
local cyclic presentation of a smooth multiple fibre,
this removes the transverse multiplicity and leaves a
translation quotient on the reduced fibre, giving exactly
the stated isogeny. Extra branch points over ordinary
fibres make the total map branched there, but its map on
the reduced elliptic fibre is an isomorphism.

On this principal bundle the source's explicit
presentation is
\[
 Z'_0=\Theta^*/\langle\tau\rangle,
 \quad |\tau|>1,\quad \deg\Theta>0,
 \qquad
 \Pic^0(Z'_0)=p'^*\Pic^0(B')\times\C^*.
 \tag{C.INT.4}
\]
In the last factor a scalar $\alpha$ is the automorphy
factor of a line bundle. Its restriction to
$E'=\C^*/\langle\tau\rangle$ is that same automorphy
line; its kernel is the discrete subgroup
$\langle\tau\rangle$. Consequently the identity
component of the restriction kernel is precisely the
compact group $p'^*\Pic^0(B')$.
Our assumed trivial restriction on $E$ makes
$\psi^*\rho(\Pic^0(Y))$ lie in that identity component.
Gauduchon degree is a continuous additive homomorphism
to $\R$, and is zero on any compact subgroup. On the
other hand this image contains
$\OO_{Z'_0}(q\psi^*A_0)$, the line of a nonzero effective
divisor, which has strictly positive Gauduchon degree.
This contradiction proves the first assertion in both
cases, and proves the stated nonconstant restriction.
Riemann--Hurwitz gives genus at most one for the
normalization of the nonconstant image of $E$, so it is
elliptic.

For the last assertion, an effective divisor of zero
real Chern class on a minimal class VII surface is,
by \cite[Remark~2.7]{DlousskyTeleman2012}, an elliptic
divisor on a Hopf surface or a positive multiple of the
rational cycle on an Enoki surface. If the source states
integral homological triviality, first multiply the
divisor by an integer killing its torsion integral
class and use Poincar\'e duality; the resulting support
and coefficient conclusion descends to the original
divisor. Formula (C.INT.2) also applies to an Enoki
minimal model. Thus the adjacent divisor on its
resolution is the full pullback of a rational cycle;
all its curve images have rational normalization.
On a Hopf or non-K\"ahler elliptic neighbouring
resolution, the adjacent divisor is the full pullback
of a nonzero divisor consisting of elliptic curves.
At least one of their strict transforms must map
into the same intersection. It cannot map to a point
or a rational-normalized curve by the result just
proved. This is the required contradiction.
\end{proof}

\begin{lemma}[C.2.INT.2; Proved exclusion of Enoki components]
\label{st:C.2.INT.2}
Let $f:Y\to B$ be a proper connected-fibre map from a
smooth compact complex threefold with $H^2(Y,\Q)=0$
and general fibre a complex two-torus. No component
of a reducible fibre has normalization resolution
with Enoki minimal model.
\end{lemma}
\begin{proof}
Write the full fibre as $F=\sum_i m_iS_i$. Its
normalization resolutions $Z_i$ are non-K\"ahler by
C.3.4. A neighbouring component gives on each $Z_i$
a nonzero effective divisor $E_{ij}=h_i^*S_j$ of
zero real Chern class. On a class VII minimal model
the precise classification in C.2.INT.1 makes that
model Hopf or Enoki. The other non-K\"ahler minimal
models are elliptic surfaces, including both types
of Kodaira surfaces. C.2.INT.1 rules out an interface
between either such type and an Enoki component.
Thus connectedness of the fibre shows that, if an
Enoki component occurs, every component has Enoki
minimal model.

Let $D_i^{\rm self}\ge0$ be the effective resolution
conductor for $S_i$ from A.5.3, and set
$D_i=D_i^{\rm self}+\sum_{j\ne i}E_{ij}$. This is
also the resolution conductor for the total reduced
support. Adjunction specifies the exact comparisons
\[
 K_{Z_i}+D_i=h_i^*(K_Y+\sum_j S_j),\qquad
 K_{Z_i}+D_i^{\rm self}=h_i^*(K_Y+S_i).
 \tag{C.INT.5}
\]
The effective numerically anticanonical divisor on
an Enoki minimal model makes it parabolic Inoue;
this is the Enoki case in
\cite[Lemma~2.22 and the following paragraph]{Dloussky2006}.
Write $E_i^0$ for its elliptic curve and $C_i$ for
its reduced rational cycle. They are disjoint and
\[
 K_{Z_i^0}+E_i^0+C_i=0.
 \tag{C.INT.6}
\]
The pushforward of $D_i$ has the form
$E_i^0+k_i' C_i$, with $k_i'\ge0$, by intersection
with $E_i^0$ and the kernel of the cycle intersection
matrix. Each nonzero adjacent divisor has the exact
form $E_{ij}=n_{ij}\sigma_i^*C_i$, $n_{ij}\ge1$.
In particular $k_i'\ge\sum_{j\ne i}n_{ij}\ge1$.

The canonical evaluation used in A.4.4 gives, on a
neighbourhood of this fibre,
$K_Y=\OO_Y(\sum_i k_iS_i)$ for integers $k_i\ge0$;
the line from the base is trivial on the chosen disc.
Put $r_i=(k_i+1)/m_i$. The principal full-fibre
relation yields the following equality in
$\Pic(Z_i)\otimes\Q$:
\[
 K_{Z_i}+D_i
   =\sum_{j\ne i}m_j(r_j-r_i)E_{ij}.
 \tag{C.INT.7}
\]
Indeed subtract $r_i h_i^*F=0$ from the first
formula (C.INT.5). Choose $i$ for which $r_i$ is
maximal. The negative of the right side of
(C.INT.7) is an effective rational divisor $A_i$.
Thus $K_{Z_i}+D_i+A_i$ is torsion after clearing
denominators. Push to its minimal model and use
(C.INT.6). Taking Gauduchon degree gives
\[
 k_i'+\sum_{j\ne i}m_j(r_i-r_j)n_{ij}=1,
 \tag{C.INT.8}
\]
since $\deg\OO(C_i)>0$. All terms after $k_i'$
are nonnegative, and $k_i'\ge\sum n_{ij}\ge1$.
Equality forces $k_i'=1$, exactly one neighbour,
$n_{ij}=1$, and $r_j=r_i$. That neighbour is again
a maximum. Propagating along the connected
intersection graph proves that there are exactly
two components $S_1,S_2$, all ratios are equal,
and on both minimal models
\[
 (D_i^{\rm self})_0=E_i^0,
 \qquad (E_{ij})_0=C_i.
 \tag{C.INT.9}
\]

We retain the geometric information in the actual
blowups and normalization maps. Every topologically
torsion line descends through a point blowup, by the
integral Picard blowup formula. Applying the
canonical divisor formula to the second equality
(C.INT.5) therefore gives at every blowup
\[
 (D_i^{\rm self})_{\rm new}
   =\sigma^*(D_i^{\rm self})_{\rm old}-E_{\rm exc}.
 \tag{C.INT.10}
\]
Starting with the smooth multiplicity-one elliptic
curve $E_i^0$, effectivity permits centres only on
its current strict transform, and the new
exceptional coefficient is zero. Hence
$D_i^{\rm self}$ is that elliptic strict transform,
and all these centres are disjoint from $C_i$.
The external cycle $\sigma_i^*C_i=C_i$ is therefore
unchanged, and disjoint from the self conductor.

It follows that $S_i$ is normal near the external
intersection. Here is the needed resolution
comparison at possible exceptional points. Any
resolution-exceptional curve outside $D_i$ is
disjoint from it: adjunction and minimality give
$K_{Z_i}\cdot E\ge0$, while $D_i\cdot E\ge0$,
so $(K_{Z_i}+D_i)\cdot E=0$ forces $D_i\cdot E=0$.
Connected resolution fibres cannot join the
elliptic self conductor to the external rational
cycle. The conductor on the normal surface is a
reflexive ideal, by the duality calculation in
A.5.3, so it has no isolated additional support
on that cycle. Thus normalization is an isomorphism
near its image. Resolution over this image is
crepant by the second formula (C.INT.5), since
$D_i^{\rm self}$ is disjoint from it. Its
exceptional curves are consequently $(-2)$
curves in $C_i$. A whole cycle cannot be contracted,
because its intersection matrix is not negative
definite. The connected fibres which are
contracted are proper chains in that cycle.

The intersection scheme $D=S_1\cap S_2$ is a
complete-intersection curve in $Y$. Its pullback
has reduced divisor $C_i$, so it is generically
reduced; its Cohen--Macaulay property then makes
it reduced everywhere. It is connected, since it
is the image of $C_i$. The equality of all $r_i$
means $K_Y+S_1+S_2=rF$ in $\Pic(Y)\otimes\Q$
on the fibre neighbourhood. Adjunction gives
$\omega_D$ torsion, hence $p_a(D)=1$.
Topologically the map $C_i\to D$ contracts just
the proper chains described above and is an
isomorphism elsewhere. It preserves the graph
circle and all its remaining components are
rational. Thus its normalized components are
$\PP^1$, its graph has exactly one cycle, and
$\sum\delta_x=\sum(r_x-1)$. Each contracted
chain has only two ends; every singular point
therefore has two branches and $\delta=1$.
The colength-one ring comparison (C.VII.6)
proves that these are ordinary nodes. At least
one node occurs.

At such a node the complete-intersection curve
has embedding dimension two. At least one of
the two surface equations has a nonzero linear
part; take it to be $z$. On $z=0$ the other
equation has an ordinary node. The holomorphic
Morse lemma with parameter $z$ therefore writes
it as $xy+h(z)$. Normality of that surface near
$D$ excludes $h\equiv0$, and a change of the
$z$ coordinate writes $h(z)=z^n$, $n\ge1$.
Consequently the original fibration has local
equation
\[
 f=u(x,y,z)\,z^a(xy+z^n)^b,
 \qquad a,b,n\ge1,
 \tag{C.INT.11}
\]
where $u$ is a unit and $a,b$ are the two full
fibre multiplicities. The unit can be removed
by an actual biholomorphic change of variables.
Put $d=2(a+nb)$ and choose a local logarithm of
$u$. For $f_0=z^a(xy+z^n)^b$, weighted homogeneity
gives
\[
 \Phi(x,y,z)=(x u^{n/d},y u^{n/d},z u^{2/d}),
 \qquad f_0\circ\Phi=u f_0.
 \tag{C.INT.12}
\]
The derivative of $\Phi$ at the origin is
invertible, proving the required equivalence
of the actual local fibres.

For $t\ne0$ small, take a simply connected
sector in the punctured $z$-plane and one
branch of $(t z^{-a})^{1/b}$. An open part
of the actual smooth fibre is
\[
 xy=g(z):=(t z^{-a})^{1/b}-z^n.
 \tag{C.INT.13}
\]
Consecutive zeros on this logarithmic branch
have angular separation $2\pi b/(a+nb)<2\pi$;
two such zeros fit in one sector. They are
simple, since at either one
$g'(z)=-(a/b+n)z^{n-1}\ne0$.
Join them by a smooth embedded path with no
other zero and take the circle $|x|=|y|$
over its interior. Explicitly, for a continuous
choice of $\arg g$ along that interior, it is
\[
 x=|g|^{1/2}e^{i(\arg g/2+\theta)},\qquad
 y=|g|^{1/2}e^{i(\arg g/2-\theta)}.
 \tag{C.INT.14}
\]
The circle collapses to a smooth disc at each
simple zero, giving an embedded sphere.
The two $x,y$ terms of the standard K\"ahler
form cancel on this sphere, and the $z$ term
vanishes on the one-dimensional path. The
sphere is Lagrangian; its normal bundle is
$JTS^2$ with the opposite Euler orientation,
so its self-intersection is $-2$.
The whole sphere lies in the local chart:
$|z|=O(|t|^{1/(a+nb)})$ and
$|x|,|y|=O(|t|^{n/(2(a+nb))})$.
Its image in the compact smooth torus fibre
would be null-homotopic, because a torus has
contractible universal cover. Its integral
homology class would therefore be zero and
its self-intersection zero, contradicting
$-2$. This contradiction excludes the Enoki
alternative.
\end{proof}

\begin{theorem}[C.2.INT.3; Proved sign and equality for every reducible fibre]
\label{st:C.2.INT.3}
Let $f:Y\to B$ satisfy the hypotheses of C.2.INT.2,
and let $S$ be a connected reducible fibre support.
Then
\[
 e(S)\ge0.
 \tag{C.INT.15}
\]
Equality holds precisely when the normalization of $S$
is a disjoint union of smooth minimal Hopf or
non-K\"ahler elliptic surfaces, and the reduced
conductor image and its reduced inverse image are
disjoint unions of smooth elliptic curves. The maps
between these elliptic curves are finite unramified
covers. This criterion allows arbitrary conductor
scheme structure and arbitrary unequal multiplicities
of the full fibre.
\end{theorem}
\begin{proof}
Let $\nu:T=\coprod T_i\to S$ be normalization,
$\pi_i:Z_i\to T_i$ its minimal resolutions, and
$\sigma_i:Z_i\to Z_i^0$ the minimal-model maps.
By C.2.INT.1--C.2.INT.2 the $Z_i^0$ are Hopf or
non-K\"ahler elliptic surfaces. They have no rational
curves and have both Euler characteristics
$e(Z_i^0)=\chi(\OO_{Z_i^0})=0$.
Every rational curve of $Z_i$ is exceptional for
$\sigma_i$; their total number is exactly the number
$N_i$ of point blowups, and $e(Z_i)=N_i$.
The elliptic strict transforms cannot be contracted
by $\pi_i$, by C.2.INT.1. Thus every
$\pi_i$-exceptional curve is $\sigma_i$-exceptional.

The connected fibres of $\pi_i$ consequently map to
points under $\sigma_i$. Normality and
$(\pi_i)_*\OO_{Z_i}=\OO_{T_i}$ descend the latter
map to a proper birational morphism
$\tau_i:T_i\to Z_i^0$, with
\[
 \sigma_i=\tau_i\pi_i.
 \tag{C.INT.16}
\]
This is an actual factorization, not a replacement
of the normalization by its minimal model.
Since point blowups have $R\sigma_{i*}\OO_{Z_i}
=\OO_{Z_i^0}$, the low-degree Leray sequence gives
\[
 0\longrightarrow R^1\tau_{i*}\OO_{T_i}
 \longrightarrow 0
 \longrightarrow \tau_{i*}R^1\pi_{i*}\OO_{Z_i}
 \longrightarrow R^2\tau_{i*}\OO_{T_i}=0.
 \tag{C.INT.17}
\]
The last vanishing follows from fibre dimension at
most one. The middle sheaf $R^1\pi_{i*}\OO_{Z_i}$
is supported at finitely many points, and its
pushforward under $\tau_i$ detects it. Hence it is
zero. This proves that $T_i$ has rational
singularities, and in particular
$\chi(\OO_{T_i})=\chi(\OO_{Z_i})=0$.

Let $\Gamma$ be the full conductor image scheme,
$A_i\subset T_i$ its full inverse-image conductor,
and let $D_i\ge0$ be the resolution conductor from
A.5.3. Its line relation
$K_{Z_i}+D_i=h_i^*(K_Y+S)$ has torsion topological
first Chern class. That line descends through
$\sigma_i$. Consequently the exact canonical
divisor formula gives
\[
 D_i=\sigma_i^*D_i^0-K_{Z_i/Z_i^0},
 \qquad D_i^0=(\sigma_i)_*D_i.
 \tag{C.INT.18}
\]
The divisor $D_i^0$ is nonzero, since $D_i$ contains
a nonzero adjacent component divisor whose
minimal-model pushforward is nonzero by (C.INT.2).
It is a sum of positive multiples of disjoint
smooth elliptic curves. At each point blowup,
(C.INT.18) has coefficient
$\operatorname{mult}_pD_{\rm old}-1$ on the new
exceptional curve. Effectivity makes the centre
lie on the old support. Starting from the disjoint
elliptic curves, the reduced support consequently
acquires only rational leaves and subdivisions of
edges. Each of its connected components is one
elliptic curve with rational trees attached.
Writing $r_i$ for its number of rational curves,
we obtain
\[
 e((D_i)_{\rm red})=r_i.
 \tag{C.INT.19}
\]

We next prove a simultaneous cohomological
comparison on the entire, possibly nonreduced,
conductor. Relative duality applied to
$R\pi_{i*}\OO_{Z_i}=\OO_{T_i}$ gives
$R\pi_{i*}\omega_{Z_i}=\omega_{T_i}$.
Finite conductor duality identifies
$\omega_{T_i}\otimes\nu^*\omega_S^{-1}$ with
the conductor ideal $\mathcal I_{A_i/T_i}$.
The divisor sequence on $Z_i$, using
$\OO_{Z_i}(-D_i)=\omega_{Z_i}\otimes h_i^*\omega_S^{-1}$,
therefore gives
\[
 R\pi_{i*}\OO_{D_i}=\OO_{A_i}.
 \tag{C.INT.20}
\]
The identical calculation for $\sigma_i$, using
(C.INT.18), gives
\[
 R\sigma_{i*}\OO_{D_i}=\OO_{D_i^0},
 \qquad R\tau_{i*}\OO_{A_i}=\OO_{D_i^0}.
 \tag{C.INT.21}
\]
For completeness the canonical pushforward for a
point blowup is elementary: its relative canonical
line is $\OO(E)$, and
$0\to\OO\to\OO(E)\to\OO_E(-1)\to0$
has zero cohomology in its last term. It gives
$R\sigma_*\omega=\omega$; composition proves the
version used in (C.INT.21).

For $M\in\Pic^0(Y)$ its pullback to $Z_i$ descends
to a line $M_i^0\in\Pic^0(Z_i^0)$, by the Picard
blowup isomorphism. Projection formula and
(C.INT.20)--(C.INT.21) identify the actual spaces
\[
 H^0(A_i,\nu^*M|_{A_i})
   =H^0(D_i^0,M_i^0|_{D_i^0}).
 \tag{C.INT.22}
\]
For a component $nE$ of $D_i^0$, the filtration by
its powers of the defining ideal has graded lines
$M_i^0|_E\otimes\OO_E(-jE)$, $0\le j<n$.
They have degree zero, since $E^2=0$ on the
minimal model. C.2.INT.1 proves that the map from
$\Pic^0(Y)$ to $\Pic^0(E)$ is nonconstant. We can
therefore choose one $M$ avoiding, for each of
the finitely many such curves and integers $j$,
the prescribed value $\OO_E(jE)$.
Each excluded locus is a proper closed analytic
subset of the connected complex Lie group
$\Pic^0(Y)$; their finite union does not fill it.
All the graded lines are nontrivial of degree
zero on an elliptic curve and have no sections.
The filtration and (C.INT.22) prove
\[
 H^0(A_i,\nu^*M)=0\quad\hbox{for every }i.
 \tag{C.INT.23}
\]

The conductor inclusion is
$\OO_\Gamma\hookrightarrow\nu_*\OO_A$, where
$A=\coprod A_i$. It now gives $H^0(\Gamma,M)=0$.
The normalization sequence and the preceding
rational-singularity computation give
$\chi(\mathcal Q)=0$ for
$\mathcal Q=\nu_*\OO_T/\OO_S$: here
$\chi(\OO_S)=0$ is the Cartier-divisor
Riemann--Roch calculation under $H^2(Y,\Q)=0$.
Conductor duality is
$\mathcal Q=\omega_\Gamma\otimes\omega_S^{-1}|_\Gamma$.
Twisting on this compact Cohen--Macaulay curve by
an ambient topologically torsion line preserves
Euler characteristic; its degree on every
component is zero. Serre duality thus gives
\[
 \chi(\OO_\Gamma)=-\chi(\mathcal Q)=0.
 \tag{C.INT.24}
\]
The same holds after twisting by $M$, so
$H^0(\Gamma,M)=0$ implies $H^1(\Gamma,M)=0$.
Restriction to $C=\Gamma_{\rm red}$ surjects
on $H^1$, since its kernel is supported in
dimension at most one and has no $H^2$.
Consequently
\[
 \chi(\OO_C)=\chi(M|_C)=h^0(C,M|_C)\ge0.
 \tag{C.INT.25}
\]
This is where the full conductor scheme is used;
no passage to its reduction is assumed harmless.

Each normalized component of $C$ has genus zero
or one, since it is dominated by a curve of
$(D_i)_{\rm red}$. Let $n_0$ be the number of
rational-normalized components, and let
$\delta_x$ and $s_x$ be the delta invariant and
branch number of a singular point of $C$.
The exact curve normalization formulas give
\[
 e(C)=\chi(\OO_C)+n_0+
                \sum_x\bigl(\delta_x-(s_x-1)\bigr)\ge0.
 \tag{C.INT.26}
\]
Every summand is nonnegative, because
$\delta_x\ge s_x-1$. Equality holds precisely
when $C$ is a disjoint union of smooth elliptic
curves. Indeed equality first forces $n_0=0$;
then $\chi(\OO_C)=-\sum\delta_x\ge0$
forces every $\delta_x=0$. The converse is immediate.

We finish the topological comparison on each
normalization component. A $\pi_i$-exceptional
curve outside $(D_i)_{\rm red}$ is a $(-2)$ curve
disjoint from that support, by the adjunction
argument following (C.INT.10). Each connected
configuration of such curves is a tree: a cycle
or a repeated intersection would contradict
negative definiteness of its matrix with
all diagonal entries $-2$. Such fibres lie
off $A_i$, because a fibre over its conductor
meets the strict conductor and is connected.
Let $a_i$ count their irreducible curves, and
let $b_i$ count exceptional fibres contained
in $(D_i)_{\rm red}$ whose image lies outside
$A_i$. Removing the conductor and contracting
these specified fibres gives the exact equality
\[
 e(T_i\setminus(A_i)_{\rm red})
       =N_i-r_i-a_i+b_i\ge0.
 \tag{C.INT.27}
\]
A tree of $a$ rational curves has Euler number
$a+1$, so its contraction removes $a$; a
positive-conductor fibre removed upstairs but
represented by a point downstairs adds one.
These are all changes, since normalization is
an isomorphism off the conductor. Nonnegativity
follows because the $r_i$ support curves and
the $a_i$ disjoint exceptional curves are
distinct rational curves among the $N_i$
rational curves of $Z_i$.
The normalization is an isomorphism
$T\setminus A\simeq S\setminus C$. Additivity
with (C.INT.26)--(C.INT.27) proves (C.INT.15).

If $e(S)=0$, every nonnegative summand in these
formulas is zero. In particular $C$ is a
union of smooth elliptic curves, $b_i=0$,
and $r_i+a_i=N_i$. Every rational curve in
$(D_i)_{\rm red}$ must then be $\pi_i$-exceptional:
a rational curve cannot map nonconstantly to
one of these elliptic curves. The $a_i$ curves
are exceptional by definition. Thus all $N_i$
rational curves of $Z_i$ are $\pi_i$-exceptional.
If $N_i>0$, the last point blowup in
$\sigma_i$ supplies a $(-1)$ curve on $Z_i$.
Its being $\pi_i$-exceptional contradicts the
minimality of that resolution. Hence $N_i=0$.
There are now no rational curves to contract,
and C.2.INT.1 excludes contracting elliptic
curves. Therefore $Z_i=T_i=Z_i^0$ is smooth
and minimal. Its reduced conductor is a
disjoint union of elliptic curves. A finite
map between smooth elliptic curves is
unramified by Riemann--Hurwitz, proving the
stated map comparison.
Conversely these normalization and conductor
conditions give $e(T_i)=e((A_i)_{\rm red})=e(C)=0$;
the normalization formula gives $e(S)=0$.
The underlying spaces, and all comparisons
just used, are independent of the positive
full-fibre multiplicities. This completes the
proof including the equality criterion.
\end{proof}
